\documentclass[12pt]{article}

\usepackage[a4paper,margin=1in]{geometry}
\usepackage{amsmath,amssymb}
\usepackage{newtxtext,newtxmath} % Times-style font
\usepackage{fancyhdr}
\usepackage{listings}
\usepackage{xcolor}
\usepackage{graphicx}
\usepackage{float}
\usepackage{algorithm}
\usepackage{algpseudocode}
\usepackage{booktabs}

\setlength{\textfloatsep}{10pt}
\setlength{\floatsep}{10pt}

\setlength{\parindent}{0pt}
\setlength{\parskip}{5pt}

% Keep entire document at 12 pt
\AtBeginDocument{
    \fontsize{12pt}{14pt}\selectfont
}

\lstdefinestyle{code}{
    language=C,
    basicstyle=\ttfamily\fontsize{9.5pt}{11.5pt}\selectfont,
    breaklines=true,
    breakatwhitespace=false,
    frame=single,
    numbers=left,
    numberstyle=\tiny\color{gray},
    keywordstyle=\color{blue},
    commentstyle=\color{gray},
    stringstyle=\color{green!50!black},
    columns=fullflexible,
    keepspaces=true,
    showstringspaces=false
}

\lstdefinestyle{output}{
    basicstyle=\ttfamily\fontsize{9.5pt}{11.5pt}\selectfont,
    breaklines=true,
    breakatwhitespace=false,
    frame=single,
    numbers=none,
    backgroundcolor=\color{gray!6},
    columns=fullflexible,
    keepspaces=true,
    showstringspaces=false
}

% ------------------------------------------------
% Header
% ------------------------------------------------

\pagestyle{fancy}
\fancyhf{}

\fancyhead[L]{Name: Kishan Sahu{\hspace{4cm}}}
\fancyhead[R]{Enrollment No.: P26DS017{\hspace{3.5cm}}}

\renewcommand{\headrulewidth}{0.4pt}
\setlength{\headheight}{15pt}

\begin{document}

% ------------------------------------------------
% TITLE
% ------------------------------------------------

\begin{center}
\textbf{Computer Science and Engineering Department, SVNIT, Surat}\\
\textbf{M.Tech. I -- Semester I}\\
\textbf{High Performance Computing (CSDS125)}\\
\textbf{Lab Assignment 1}\\
\textbf{Dynamic Memory Allocation: Array Operations, Matrix Multiplication, Vector Dot Product, and Merge Sort}
\end{center}
\vspace{-4pt}
\hrule
\vspace{8pt}

% ==============================================================================
% Problem 1: Minimum and Maximum in Array
% ==============================================================================
\section*{Problem 1: Minimum and Maximum Element in Array using Dynamic Memory}

\subsection*{1.1 Objective}
Implement a C program to allocate memory dynamically for an array of size $n$, read the elements from the user, and find the minimum and maximum elements in a single traversal.

\subsection*{1.2 Source Code (\texttt{min\_max.c})}
\begin{lstlisting}[style=code]
#include <stdio.h>
#include <stdlib.h>

int main() {
    int n;
    printf("Enter number of elements: ");
    if (scanf("%d", &n) != 1 || n <= 0) {
        printf("Invalid size\n");
        return 1;
    }

    int *arr = (int *)malloc(n * sizeof(int));
    if (arr == NULL) {
        printf("Memory allocation failed\n");
        return 1;
    }

    printf("Enter %d elements:\n", n);
    for (int i = 0; i < n; i++) {
        scanf("%d", &arr[i]);
    }

    int min = arr[0];
    int max = arr[0];

    for (int i = 1; i < n; i++) {
        if (arr[i] < min) {
            min = arr[i];
        }
        if (arr[i] > max) {
            max = arr[i];
        }
    }

    printf("Minimum element: %d\n", min);
    printf("Maximum element: %d\n", max);

    free(arr);
    return 0;
}
\end{lstlisting}

\subsection*{1.3 Execution Output}
\begin{lstlisting}[style=output]
Enter number of elements: 5
Enter 5 elements:
10 -5 30 2 0
Minimum element: -5
Maximum element: 30
\end{lstlisting}

\subsection*{1.4 Analysis}
\begin{itemize}
    \item \textbf{Time Complexity:} $\mathcal{O}(n)$. Memory allocation takes $\mathcal{O}(1)$ time. Array population and the subsequent single-pass scan iterate from index $1$ to $n-1$, performing exactly $2(n-1)$ comparisons in the worst case.
    \item \textbf{Space Complexity:} $\mathcal{O}(n)$ auxiliary space allocated on the heap for storing $n$ integers. The algorithm maintains only $\mathcal{O}(1)$ extra variables (\texttt{min}, \texttt{max}, \texttt{i}, \texttt{n}).
    \item \textbf{Memory Management:} Explicit heap memory is allocated using \texttt{malloc} and deallocated using \texttt{free(arr)} before program termination, preventing memory leaks.
\end{itemize}

\vspace{10pt}
\hrule
\vspace{10pt}

% ==============================================================================
% Problem 2: Matrix Multiplication
% ==============================================================================
\section*{Problem 2: Matrix Multiplication using Dynamic Memory Allocation}

\subsection*{2.1 Objective}
Implement matrix multiplication $C = A \times B$ for matrices of arbitrary dimensions ($r_1 \times c_1$ and $r_2 \times c_2$) using dynamically allocated 2D pointer arrays, with dimension compatibility validation ($c_1 = r_2$).

\subsection*{2.2 Source Code (\texttt{matrix\_mult.c})}
\begin{lstlisting}[style=code]
#include <stdio.h>
#include <stdlib.h>

int main() {
    int r1, c1, r2, c2;

    printf("Enter rows and columns of Matrix A: ");
    if (scanf("%d %d", &r1, &c1) != 2 || r1 <= 0 || c1 <= 0) {
        printf("Invalid dimensions\n");
        return 1;
    }

    printf("Enter rows and columns of Matrix B: ");
    if (scanf("%d %d", &r2, &c2) != 2 || r2 <= 0 || c2 <= 0) {
        printf("Invalid dimensions\n");
        return 1;
    }

    if (c1 != r2) {
        printf("Multiplication not possible: columns of A must equal rows of B\n");
        return 1;
    }

    // Allocate Matrix A
    int **A = (int **)malloc(r1 * sizeof(int *));
    for (int i = 0; i < r1; i++) {
        A[i] = (int *)malloc(c1 * sizeof(int));
    }

    // Allocate Matrix B
    int **B = (int **)malloc(r2 * sizeof(int *));
    for (int i = 0; i < r2; i++) {
        B[i] = (int *)malloc(c2 * sizeof(int));
    }

    // Allocate Result Matrix C
    int **C = (int **)malloc(r1 * sizeof(int *));
    for (int i = 0; i < r1; i++) {
        C[i] = (int *)malloc(c2 * sizeof(int));
    }

    printf("Enter elements of Matrix A:\n");
    for (int i = 0; i < r1; i++) {
        for (int j = 0; j < c1; j++) {
            scanf("%d", &A[i][j]);
        }
    }

    printf("Enter elements of Matrix B:\n");
    for (int i = 0; i < r2; i++) {
        for (int j = 0; j < c2; j++) {
            scanf("%d", &B[i][j]);
        }
    }

    // Multiply matrices
    for (int i = 0; i < r1; i++) {
        for (int j = 0; j < c2; j++) {
            C[i][j] = 0;
            for (int k = 0; k < c1; k++) {
                C[i][j] += A[i][k] * B[k][j];
            }
        }
    }

    printf("Resultant Matrix:\n");
    for (int i = 0; i < r1; i++) {
        for (int j = 0; j < c2; j++) {
            printf("%d ", C[i][j]);
        }
        printf("\n");
    }

    // Free memory
    for (int i = 0; i < r1; i++) free(A[i]);
    free(A);

    for (int i = 0; i < r2; i++) free(B[i]);
    free(B);

    for (int i = 0; i < r1; i++) free(C[i]);
    free(C);

    return 0;
}
\end{lstlisting}

\subsection*{2.3 Execution Output}
\begin{lstlisting}[style=output]
Enter rows and columns of Matrix A: 2 3
Enter rows and columns of Matrix B: 3 2
Enter elements of Matrix A:
1 2 3
4 5 6
Enter elements of Matrix B:
7 8
9 1
2 3
Resultant Matrix:
31 19 
85 55 
\end{lstlisting}

\subsection*{2.4 Analysis}
\begin{itemize}
    \item \textbf{Time Complexity:} $\mathcal{O}(r_1 \cdot c_1 \cdot c_2)$. Matrix multiplication uses three nested loops where the outer two iterate over the rows ($r_1$) and columns ($c_2$) of the output matrix, and the innermost loop performs $c_1$ multiply-accumulate operations. For square matrices of size $n \times n$, this corresponds to standard $\mathcal{O}(n^3)$ time complexity.
    \item \textbf{Space Complexity:} $\mathcal{O}(r_1 c_1 + r_2 c_2 + r_1 c_2)$ heap memory to store matrices $A$, $B$, and $C$. In addition, $\mathcal{O}(r_1 + r_2)$ pointers are allocated for the row index structures.
    \item \textbf{Memory Management:} 2D arrays are implemented using an array of row pointers (\texttt{int**}). Each row is dynamically allocated and freed in reverse order before the row pointer table is released, guaranteeing complete deallocation without memory leakage.
\end{itemize}

\vspace{10pt}
\hrule
\vspace{10pt}

% ==============================================================================
% Problem 3: Vector Dot Product
% ==============================================================================
\section*{Problem 3: Dot Product of Two Vectors using Dynamic Memory Allocation}

\subsection*{3.1 Objective}
Compute the scalar dot product $\vec{a} \cdot \vec{b} = \sum_{i=0}^{n-1} a_i b_i$ of two vectors of user-specified size $n$ using dynamically allocated buffers.

\subsection*{3.2 Source Code (\texttt{dot\_product.c})}
\begin{lstlisting}[style=code]
#include <stdio.h>
#include <stdlib.h>

int main() {
    int n;
    printf("Enter size of vectors: ");
    if (scanf("%d", &n) != 1 || n <= 0) {
        printf("Invalid size\n");
        return 1;
    }

    int *a = (int *)malloc(n * sizeof(int));
    int *b = (int *)malloc(n * sizeof(int));

    if (a == NULL || b == NULL) {
        printf("Memory allocation failed\n");
        return 1;
    }

    printf("Enter elements of first vector:\n");
    for (int i = 0; i < n; i++) {
        scanf("%d", &a[i]);
    }

    printf("Enter elements of second vector:\n");
    for (int i = 0; i < n; i++) {
        scanf("%d", &b[i]);
    }

    long long dot_product = 0;
    for (int i = 0; i < n; i++) {
        dot_product += (long long)a[i] * b[i];
    }

    printf("Dot Product: %lld\n", dot_product);

    free(a);
    free(b);
    return 0;
}
\end{lstlisting}

\subsection*{3.3 Execution Output}
\begin{lstlisting}[style=output]
Enter size of vectors: 3
Enter elements of first vector:
1 2 3
Enter elements of second vector:
4 5 6
Dot Product: 32
\end{lstlisting}

\subsection*{3.4 Analysis}
\begin{itemize}
    \item \textbf{Time Complexity:} $\mathcal{O}(n)$. Reading the input vectors requires $2n$ operations and the linear accumulation loop requires $n$ multiplication and addition operations.
    \item \textbf{Space Complexity:} $\mathcal{O}(n)$ dynamic memory for the two vectors ($2n \times \text{sizeof(int)}$ bytes). Auxiliary stack space is $\mathcal{O}(1)$.
    \item \textbf{Numerical Robustness:} The dot product accumulator is declared as \texttt{long long} (64-bit integer) to prevent potential arithmetic overflow when accumulating products of large 32-bit integers.
\end{itemize}

\vspace{10pt}
\hrule
\vspace{10pt}

% ==============================================================================
% Problem 4: Divide and Conquer Sorting (Merge Sort)
% ==============================================================================
\section*{Problem 4: Array Sorting using Divide-and-Conquer (Merge Sort)}

\subsection*{4.1 Objective}
Implement the divide-and-conquer Merge Sort algorithm to sort an array of $n$ elements in non-decreasing order using dynamic memory allocation for the input array and merge buffers.

\subsection*{4.2 Source Code (\texttt{merge\_sort.c})}
\begin{lstlisting}[style=code]
#include <stdio.h>
#include <stdlib.h>

void merge(int arr[], int l, int m, int r) {
    int n1 = m - l + 1;
    int n2 = r - m;

    int *L = (int *)malloc(n1 * sizeof(int));
    int *R = (int *)malloc(n2 * sizeof(int));

    for (int i = 0; i < n1; i++) {
        L[i] = arr[l + i];
    }
    for (int j = 0; j < n2; j++) {
        R[j] = arr[m + 1 + j];
    }

    int i = 0, j = 0, k = l;
    while (i < n1 && j < n2) {
        if (L[i] <= R[j]) {
            arr[k++] = L[i++];
        } else {
            arr[k++] = R[j++];
        }
    }

    while (i < n1) {
        arr[k++] = L[i++];
    }
    while (j < n2) {
        arr[k++] = R[j++];
    }

    free(L);
    free(R);
}

void mergeSort(int arr[], int l, int r) {
    if (l < r) {
        int m = l + (r - l) / 2;
        mergeSort(arr, l, m);
        mergeSort(arr, m + 1, r);
        merge(arr, l, m, r);
    }
}

int main() {
    int n;
    printf("Enter number of elements: ");
    if (scanf("%d", &n) != 1 || n <= 0) {
        printf("Invalid size\n");
        return 1;
    }

    int *arr = (int *)malloc(n * sizeof(int));
    if (arr == NULL) {
        printf("Memory allocation failed\n");
        return 1;
    }

    printf("Enter %d elements:\n", n);
    for (int i = 0; i < n; i++) {
        scanf("%d", &arr[i]);
    }

    mergeSort(arr, 0, n - 1);

    printf("Sorted array:\n");
    for (int i = 0; i < n; i++) {
        printf("%d ", arr[i]);
    }
    printf("\n");

    free(arr);
    return 0;
}
\end{lstlisting}

\subsection*{4.3 Execution Output}
\begin{lstlisting}[style=output]
Enter number of elements: 7
Enter 7 elements:
38 27 43 3 9 82 10
Sorted array:
3 9 10 27 38 43 82 
\end{lstlisting}

\subsection*{4.4 Analysis}
\begin{itemize}
    \item \textbf{Algorithmic Paradigm:} Divide-and-Conquer. The problem is partitioned recursively into two halves ($T(n/2)$), sorted, and merged in linear time ($\mathcal{O}(n)$).
    \item \textbf{Recurrence Relation:} $T(n) = 2T(n/2) + \mathcal{O}(n)$, which yields a tight bound of $\Theta(n \log n)$ by the Master Theorem across best, average, and worst cases.
    \item \textbf{Space Complexity:} $\mathcal{O}(n)$ auxiliary heap space allocated dynamically during the merge step for subarray buffers $L$ and $R$, alongside $\mathcal{O}(\log n)$ recursion stack depth.
    \item \textbf{Stability:} Merge Sort is stable, meaning elements with equal keys maintain their relative original input order (achieved via \texttt{L[i] <= R[j]}).
\end{itemize}

\end{document}
